\section{Preliminaries: Normal Forms from Subgroup Chains}\label{sec:preliminaries}

The library's normal forms all follow one classical recipe from
computational group theory.  This section recalls it in mathematical
vocabulary and fixes the intuitions --- transversals as digits, towers as
positional number systems --- that the Agda modules of
\S\ref{sec:design} implement.  Standard references are Sims's monograph on
computation with finitely presented groups~\cite{sims1994computation} and
the Handbook of Computational Group Theory~\cite{holt2005handbook}; the
subgroup-presentation method goes back to Reidemeister and Schreier, in the
form given by Magnus, Karrass, and Solitar~\cite{magnus2004combinatorial}.

\subsection{Transversals and digits}

Let $H \le G$ be a subgroup.  A \emph{(right) transversal} is a set of
representatives for the right cosets $Hg$, one per coset; write $\bar g$
for the representative of $Hg$.  Every $g \in G$ then factors
\emph{uniquely} as
\[
  g \;=\; h \cdot \bar g, \qquad h \in H,\ \bar g \in \text{transversal},
\]
and the map $g \mapsto (h, \bar g)$ is a bijection
$G \cong H \times (H\backslash G)$ of sets (not of groups).  If one now
fixes a \emph{chain} of subgroups
\[
  1 = G_0 \;\le\; G_1 \;\le\; \dots \;\le\; G_n = G ,
\]
with a transversal $T_k$ for each step $G_{k-1} \le G_k$, iterating the
factorization expresses every $g \in G$ uniquely as
\[
  g \;=\; t_1 \, t_2 \cdots t_n , \qquad t_k \in T_k .
\]
The tuple $(t_1,\dots,t_n)$ is a \emph{normal form} for $g$: existence is
the iterated factorization, uniqueness is uniqueness of coset
representatives at each level.  When the indices $[G_k : G_{k-1}] = r_k$
are finite, $|G| = r_1 r_2 \cdots r_n$, and the normal form is literally a
\emph{mixed-radix positional numeral}: the $k$-th digit ranges over $r_k$
values, and enumerating tuples lexicographically enumerates the group.  For
the chain $S_1 < S_2 < \dots < S_n$ of symmetric groups the radices are
$2, 3, \dots, n$ and the numeral system is the factorial number system ---
which is exactly the carrier $\aid{NF}\;n = \top \times C\,1 \times \dots
\times C\,(n{-}1)$ built by the coset tower in \S\ref{sec:permutations}.  For a
group of order $8 \cdot 4$ presented as a product of cyclic groups of
orders $8$ and $4$ (the base of the Clifford+$T$ tower,
\S\ref{sec:examples}) the radices are $8$ and $4$.  This digit reading of
stabilizer-chain normal forms is the conceptual core of the
Schreier--Sims method for permutation
groups~\cite{sims1994computation,holt2005handbook}, where the $G_k$ are
point stabilizers; our chains are chosen syntactically instead, but the
arithmetic is the same.

\subsection{Words, presentations, and the word problem}

A \emph{presentation} $\langle X \mid R \rangle$ consists of a set $X$ of
generators and a relation $R$ on the free monoid $X^{*}$ --- a set of pairs
of words over $X$, the axioms.  We work with monoid presentations, so
words are finite sequences with concatenation and there are no formal
inverses.  The presented monoid is the quotient of $X^{*}$ by the least
congruence $\approx$ containing $R$; the \emph{word problem} asks to decide
$w \approx v$.  An interpretation of the generators in a monoid or group
$G$ extends, by freeness of $X^{*}$, to a unique monoid homomorphism
$X^{*} \to G$; when every axiom in $R$ holds in $G$ --- \emph{soundness}
--- this homomorphism respects $\approx$ and so descends to the quotient,
and it is this \emph{induced} homomorphism $X^{*}/{\approx} \to G$ that we
ask to be an isomorphism.  A presentation \emph{presents} $G$ (relative to
the interpretation) when it is; its injectivity is \emph{completeness} of
$R$ for $G$, its surjectivity says that the images of the generators
generate $G$.  A group presentation in the
classical sense is recovered by adding, for each generator, a generator
playing the role of its inverse together with the cancellation relations
--- in the library this is the \aid{Grouplike} condition of
\S\ref{sec:presentations}, which asks every generator to have a left
inverse up to $\approx$ and derives the group structure from it.

Given $H \le G$ with $G = \langle Y \mid S \rangle$, the
\emph{Reidemeister--Schreier} method produces a presentation of $H$: its
generators are the \emph{Schreier generators} $\bar t y \, \overline{ty}^{\,-1}$
($t$ in a transversal, $y \in Y$), and its relations are rewrites of the
relators of $G$ through the coset table.  The data one must supply is
finite and concrete --- the coset table
$\tau : \text{cosets} \times Y \to \text{cosets}$ recording $Hty = H\tau(t,y)$,
together with, for each step, the element of $H$ by which $ty$ differs from
its representative.  Classical treatments run this as an algorithm
(Todd--Coxeter coset enumeration~\cite{sims1994computation,
holt2005handbook}); we instead take a \emph{claimed} coset table
as input and \emph{verify} it, by exhibiting five equational
hypotheses relating the table, the section, and the relations
(\S\ref{sec:engine}).  Verification rather than search is the right
division of labour in a proof assistant: tables are found offline (by
computation, or on paper), and the checked artifact certifies them.
Bian and Selinger pioneered this monoid-level, verified use of
Reidemeister--Schreier in Agda for their gate-group
presentations~\cite{bian2023cliffordt2,bian2023cliffordcs3,bian2023thesis};
we turned it into a reusable module with the hypotheses packaged as
a record.

\subsection{Specialization to circuits}

For circuits the chain is not chosen --- it is handed to us by the wire
structure.  Writing $\mathsf{Cir}\,k$ for the group (or monoid) presented
by $k$-wire circuits modulo a relation set, the shift operation
$c \mapsto c{\uparrow}$ embeds
\[
  \mathsf{Cir}\,1 \;\le\; \mathsf{Cir}\,2 \;\le\; \dots \;\le\;
  \mathsf{Cir}\,k \;\le\; \mathsf{Cir}\,(k{+}1) \;\le\; \dots
\]
as ``the circuits that do not touch the bottom wire''.  A normal form for
$\mathsf{Cir}\,(k{+}1)$ over $\mathsf{Cir}\,k$ is then a transversal of
this subgroup --- for permutations, the staircases of
\S\ref{sec:sn-cosets}; for Clifford-style groups, Selinger's
staircase normal forms~\cite{selinger2015clifford} and their qutrit
analogues~\cite{li2025qutrit} --- and a normal form for the whole group is
a tower of such levels.  The library's \aid{CosetTower}
(\S\ref{sec:engine}) is precisely this specialization: an $\mathbb{N}$-indexed
family of presentations, a coset type per level, and a verified table per
level, folded into a normal form whose carrier grows by one digit per
level.  Where the subgroup chain branches instead of nesting --- two
subgroups sharing a common base, as for the Clifford+$T$ groups ---
the appropriate gluing is the \emph{amalgamated free product}, whose
classical alternating normal form~\cite{magnus2004combinatorial} we
also verify (\S\ref{sec:constructions}).
